\section*{Capítulo \ref{ch:b1:det} --- Determinantes y sistemas lineales}

\begin{solution}{exo:b1:det:1}
$3 \times 2 - 1 \times 5 = 1$.

Segundo: $L_2 \leftarrow L_2 - L_1$, $L_3 \leftarrow L_3 - L_2$
(sobre las filas originales) dan las filas $(1,2,3), (3,3,3),
(3,3,3)$: dos filas iguales, determinante $0$. (Sarrus lo confirma:
$45 + 84 + 96 - 105 - 48 - 72 = 0$.)

Tercero: es un Vandermonde con $x = 1, 2, 3$
(\cref{ex:b1:det:vandermonde}): $(2-1)(3-1)(3-2) = 2$.
\end{solution}

\begin{solution}{exo:b1:det:2}
El determinante es (súmense todas las columnas a la primera y
sáquese factor común) $(\lambda + 2)$ por
\[
\begin{vmatrix}
1 & 1 & \lambda\\ 1 & \lambda & 1\\ 1 & 1 & 1
\end{vmatrix}
= -(\lambda - 1)^2
\]
(límpiese con $L_1 \leftarrow L_1 - L_3$, $L_2 \leftarrow L_2 - L_3$
y desarróllese), lo que da $\det = -(\lambda+2)(\lambda-1)^2$. Es
base $\iff \det \neq 0 \iff \lambda \notin \{1, -2\}$.
\end{solution}

\begin{solution}{exo:b1:det:3}
Primer sistema: $\det = -7$;
$x = \frac{1}{-7}\begin{vmatrix} 5 & 1\\ 4 & -2\end{vmatrix}
= \frac{-14}{-7} = 2$,
$y = \frac{1}{-7}\begin{vmatrix} 2 & 5\\ 3 & 4\end{vmatrix}
= \frac{-7}{-7} = 1$. Comprobación: $2(2) + 1 = 5$; $3(2) - 2 = 4$.

Segundo sistema: tras $L_2 - L_1$ y $L_3 - 2L_1$, las filas pasan a
ser $(1,1,1)$, $(0,-2,0)$, $(0,-1,-3)$, luego
\[
\det A = \begin{vmatrix} 1&1&1\\ 1&-1&1\\ 2&1&-1\end{vmatrix}
= 1 \times \begin{vmatrix} -2 & 0\\ -1 & -3\end{vmatrix} = 6 .
\]
Cramer, sustituyendo columnas por $(6,2,1)^{\mathsf T}$:
\[
x = \frac{6}{6} = 1, \qquad
y = \frac{12}{6} = 2, \qquad
z = \frac{18}{6} = 3
\]
(los numeradores calculados igual). Comprobación: $1 + 2 + 3 = 6$;
$1 - 2 + 3 = 2$; $2 + 2 - 3 = 1$.
\end{solution}

\begin{solution}{exo:b1:det:4}
Redúzcase la matriz ampliada: $L_2 \leftarrow L_2 - 2L_1$,
$L_3 \leftarrow L_3 - L_1$:
\[
\begin{pmatrix}
1 & 2 & -1 & 1 & 1\\
0 & 0 & 3 & -3 & 3\\
0 & 0 & 3 & -3 & 3
\end{pmatrix}
\to
\begin{pmatrix}
1 & 2 & -1 & 1 & 1\\
0 & 0 & 1 & -1 & 1\\
0 & 0 & 0 & 0 & 0
\end{pmatrix}.
\]
Incógnitas principales $x, z$; incógnitas \omterm{def:b1:vspaces:free}{libres} $y, t$. Sustitución
hacia atrás: $z = 1 + t$, $x = 1 - 2y + z - t = 2 - 2y$. \omterm{def:b1:logic:sets}{Conjunto}
de soluciones:
\[
\{(2 - 2y,\; y,\; 1 + t,\; t) : y, t \in \R\}
= (2, 0, 1, 0) + \operatorname{Vect}\bigl((-2,1,0,0),\,
(0,0,1,1)\bigr),
\]
un plano afín (de dimensión $2 = 4 - \operatorname{rk} 2$) de $\R^4$.
\end{solution}

\begin{solution}{exo:b1:det:5}
Inducción; $n = 1$ es el producto vacío $= 1$. Para el paso, háganse
$C_k \leftarrow C_k - x_1 C_{k-1}$ para $k = n, n-1, \dots, 2$ (en
ese orden, de modo que cada operación use una columna aún no
modificada). La primera fila pasa a ser $(1, 0, \dots, 0)$; y en la
fila $i \geq 2$, la entrada $k$-ésima pasa a ser
$x_i^{k-1} - x_1 x_i^{k-2} = x_i^{k-2}(x_i - x_1)$. Desarrollando por
la primera fila y sacando factor $(x_i - x_1)$ de cada fila $i$:
\[
V(x_1, \dots, x_n)
= \prod_{i=2}^{n} (x_i - x_1)\cdot V(x_2, \dots, x_n),
\]
y la hipótesis de inducción completa el producto
$\prod_{i<j}(x_j - x_i)$.
\end{solution}

\begin{solution}{exo:b1:det:6}
Desarrollando $D_n$ por la primera fila:
$D_n = 2 D_{n-1} - 1\cdot\begin{vmatrix} 1 & \ast\\ 0 &
D_{n-2}\text{-bloque} \end{vmatrix}$; y el segundo determinante,
desarrollado por su primera columna, es $D_{n-2}$. Luego
$D_n = 2D_{n-1} - D_{n-2}$, es decir,
$D_n - D_{n-1} = D_{n-1} - D_{n-2}$: las diferencias son constantes e
iguales a $D_2 - D_1 = 1$. Por tanto, $D_n = D_1 + (n - 1) = n + 1$.
(Compruébese: $D_2 = 3$, y el caso $3\times3$ es
el~\cref{ex:b1:det:cofactor}: $D_3 = 4$.)
\end{solution}

\begin{solution}{exo:b1:det:7}
$C_1 \leftarrow C_1 + C_2 + C_3$ vuelve constante la primera columna,
igual a $(m+2)$; sáquese factor común:
\[
\det = (m+2)\begin{vmatrix}
1 & 1 & m\\ 1 & m & 1\\ 1 & 1 & 1
\end{vmatrix}
\overset{L_1 - L_3,\ L_2 - L_3}{=}
(m+2)\begin{vmatrix}
0 & 0 & m-1\\ 0 & m-1 & 0\\ 1 & 1 & 1
\end{vmatrix}
= (m+2)\cdot\bigl(-(m-1)^2\bigr)
\]
(desarróllese por la primera columna: la única entrada $1$ lleva
signo $+$, y el determinante $2 \times 2$ restante es
$0 \cdot 0 - (m-1)(m-1) = -(m-1)^2$).

Discusión. $m \notin \{1, -2\}$: sistema de Cramer; por la simetría
de las ecuaciones, $x = y = z$, y cada ecuación da $(m + 2)x = 1$:
solución única
$\bigl(\frac{1}{m+2}, \frac{1}{m+2}, \frac{1}{m+2}\bigr)$. $m = 1$:
las tres ecuaciones dicen todas $x + y + z = 1$: las soluciones
forman el plano afín $x + y + z = 1$. $m = -2$: sumando las tres
ecuaciones se obtiene $0 = 3$: \omterm{def:b1:logic:sets}{conjunto} de soluciones vacío.
\end{solution}

\begin{solution}{exo:b1:det:8}
($\Rightarrow$) Si $A^{-1}$ tiene entradas enteras:
$\det A \cdot \det A^{-1} = 1$ con los dos determinantes enteros
(sumas de productos de entradas): dos enteros de producto $1$ son
los dos $\pm1$.

($\Leftarrow$) La fórmula de los cofactores
$A^{-1} = \frac{1}{\det A}\operatorname{Com}(A)^{\mathsf T}$
(comprobada para $n \leq 3$ por desarrollo directo y admitida en
general) tiene $\operatorname{Com}(A)$ con entradas enteras (cada
cofactor es un determinante entero); y dividir por
$\det A = \pm 1$ mantiene los enteros.
\end{solution}

\begin{solution}{exo:b1:det:9}
Súmense todas las columnas a la primera: cada entrada de la nueva
primera columna es $a + nb$; sáquese factor común, de modo que la
primera columna sea de unos. Después, las \omterm{met:b1:matrices:gauss}{operaciones por filas}
$L_i \leftarrow L_i - L_1$ ($i \geq 2$) limpian todas las entradas
por debajo del $1$ superior izquierdo y dejan $a$ en la diagonal y
$0$ en el resto de esas filas: la matriz es triangular superior con
diagonal $(1, a, \dots, a)$. Por tanto,
\[
\det(aI + bJ) = (a + nb)\, a^{\,n-1} ,
\]
no nulo si y solo si $a \neq 0$ y $a + nb \neq 0$: la condición
del~\cref{exo:b1:matrices:9}.
\end{solution}

\begin{solution}{exo:b1:det:10}
Operaciones por bloques (cada una, una composición de las $n$
operaciones escalares correspondientes, permitidas por
el~\cref{thm:b1:det:props}~(1)):
\[
\begin{vmatrix} A & B\\ B & A\end{vmatrix}
\overset{C_1 \leftarrow C_1 + C_2}{=}
\begin{vmatrix} A + B & B\\ A + B & A\end{vmatrix}
\overset{L_2 \leftarrow L_2 - L_1}{=}
\begin{vmatrix} A + B & B\\ 0 & A - B\end{vmatrix}
= \det(A+B)\,\det(A-B),
\]
usando en el último paso la regla triangular por bloques.
\end{solution}

\begin{solution}{exo:b1:det:11}
$C_1 \leftarrow C_1 + C_2 + C_3$ vuelve constante la primera columna,
igual a $(a + b + c)$; sáquese factor común y háganse después
$L_2 \leftarrow L_2 - L_1$, $L_3 \leftarrow L_3 - L_1$:
\[
\Delta = (a+b+c)\begin{vmatrix}
1 & b & c\\
0 & a - b & b - c\\
0 & c - b & a - c
\end{vmatrix}
= (a+b+c)\bigl[(a-b)(a-c) + (b-c)^2\bigr],
\]
y, desarrollando,
$(a-b)(a-c) + (b-c)^2 = a^2 + b^2 + c^2 - ab - bc - ca$. Sobre $\C$,
con $j^3 = 1$ y $1 + j + j^2 = 0$:
\begin{align*}
(a + jb + j^2c)(a + j^2b + jc)
&= a^2 + b^2 + c^2 + (j + j^2)(ab + bc + ca)\\
&= a^2 + b^2 + c^2 - ab - bc - ca ,
\end{align*}
de donde la factorización completa. (Estructuralmente: la columna
$(1, j, j^2)^{\mathsf T}$ cumple
$M\,(1, j, j^2)^{\mathsf T} = (a + jb + j^2c)(1, j,
j^2)^{\mathsf T}$, y análogamente para $j^2$ y $1$: los tres factores
son los tres «valores propios» del circulante, historia que se
sistematiza en el volumen del segundo año.)
\end{solution}

\begin{solution}{exo:b1:det:12}
Escríbase $r = \operatorname{rk} A$.

\emph{Existe una submatriz $r \times r$ invertible.} Elíjanse $r$
columnas \omterm{def:b1:vspaces:free}{libres} de $A$ y sea $B \in \mathcal{M}_{n,r}$ la matriz
que forman: $\operatorname{rk} B = r$. Como el rango por filas es
igual al rango por columnas (\cref{thm:b1:matrices:rank}), $B$ tiene
$r$ filas \omterm{def:b1:vspaces:free}{libres}; y quedarse con esas filas da una submatriz
$r \times r$ de $A$ de rango $r$, es decir, invertible.

\emph{No hay ninguna mayor.} Sea $S$ una submatriz $s \times s$
invertible, tomada de las columnas $j_1, \dots, j_s$ y las filas
$i_1, \dots, i_s$ de $A$. Si se anula una combinación
$\sum_k \lambda_k C_{j_k} = 0$ de las columnas \emph{completas}
correspondientes, entonces leer solo las filas $i_1, \dots, i_s$ da
$\sum_k \lambda_k S_k = 0$ sobre las columnas de $S$, luego todos los
$\lambda_k = 0$ ($S$ es invertible): las columnas
$C_{j_1}, \dots, C_{j_s}$ de $A$ son \omterm{def:b1:vspaces:free}{libres} y
$s \leq \operatorname{rk} A = r$.

Por tanto, $\operatorname{rk} A$ es exactamente el mayor tamaño de
una submatriz invertible.
\end{solution}

\begin{solution}{pb:b1:det:1}
\textbf{1.} $V(1,2,3,4) = (2-1)(3-1)(4-1)(3-2)(4-2)(4-3) =
1 \cdot 2\cdot 3\cdot 1\cdot 2\cdot 1 = 12$. La interpolación en
nodos distintos pide los coeficientes de $P$ que resuelven $W c = y$
con $W = (x_i^{\,j-1})$, y $\det W = V \neq 0$: un sistema de Cramer.

\textbf{2.} Recórranse las columnas de izquierda a derecha. $C_1$ es
la columna constante $P_0(x_i) = 1$ ($P_0$ es \omterm{def:b1:poly:def}{mónico} de grado $0$).
Supóngase que las columnas $1, \dots, j-1$ ya se han reducido a las
potencias puras $1, x_i, \dots, x_i^{\,j-2}$. Como
$P_{j-1} = X^{j-1} + \sum_{k < j-1}\alpha_k X^k$, restar a $C_j$ la
combinación $\sum_k \alpha_k\,(\text{columna de } x_i^k)$ ---una
operación que no cambia el determinante--- deja la columna de
potencias puras $x_i^{\,j-1}$. Tras la última columna, la matriz es
la de Vandermonde: $\det = V(x_1, \dots, x_n)$.

\textbf{3.} Los \omterm{def:b1:poly:def}{polinomios} $(j-1)!\,B_{j-1}$ son \omterm{def:b1:poly:def}{mónicos} de grado
$j - 1$, así que la pregunta 2 da
\[
\det\bigl(B_{j-1}(m_i)\bigr)
= \frac{V(m_1, \dots, m_n)}{0!\,1!\cdots(n-1)!} .
\]
El miembro izquierdo es el determinante de una matriz con entradas
\emph{enteras} ($B_k$ es de valores enteros en $\Z$: preguntas 16--17
del problema del fin de semana \cref{pb:b1:vspaces:1}) y, por tanto,
un entero; y es positivo, ya que $V(m_1, \dots, m_n) > 0$ para
$m_1 < \dots < m_n$. Luego el superfactorial \omterm{def:b1:arith:divides}{divide} al producto de
todas las diferencias dos a dos.

\textbf{4.} Sáquese factor $x_i$ de cada fila $i$:
$\det(x_i^{\,j})_{j = 1..n} = x_1\cdots x_n\, \det(x_i^{\,j-1}) =
x_1\cdots x_n\,V$.

\textbf{5.} $(W^{\mathsf T}W)_{ij} = \sum_k x_k^{\,i-1}
x_k^{\,j-1} = p_{i+j-2}$: $S = W^{\mathsf T}W$. Por tanto,
$\det S = \det(W^{\mathsf T})\det W = V^2$
(\cref{thm:b1:det:props}~(2),(4)). Para $x_i$ reales:
$\det S = V^2 \geq 0$, y $S$ es invertible si y solo si $V \neq 0$,
si y solo si los $x_i$ son distintos dos a dos --- un test de signo
definido y calculable solo a partir de las sumas de potencias.

\textbf{6.} Las condiciones de interpolación
$\sum_{k} c_k\,x_i^{\,k} = y_i$ forman el sistema $Wc = y$; y
$\det W = V \neq 0$ da de golpe la existencia y la unicidad. Esta es
la tercera demostración del libro: fórmula explícita en
el~\cref{thm:b1:poly:lagrange}, argumento de núcleo en
el~\cref{ex:b1:linmaps:interpolation} y Cramer aquí.

\textbf{7.} Cramer: $c_{n-1} = \det W'/\det W$, donde $W'$ es $W$ con
su última columna sustituida por $y$. Desarrollando $\det W'$ por esa
columna:
\[
\det W' = \sum_{i=1}^n (-1)^{i+n} y_i\,V(x_1, \dots, \widehat{x_i},
\dots, x_n) .
\]
Ahora bien, $V = V(\setminus i)\cdot\prod_{j<i}(x_i -
x_j)\prod_{j>i} (x_j - x_i)$, y convertir el segundo producto cuesta
$(-1)^{n-i}$:
\[
(-1)^{i+n}\,\frac{V(\setminus i)}{V}
= \frac{(-1)^{i+n}(-1)^{n-i}}{\prod_{j\neq i}(x_i - x_j)}
= \frac{1}{\prod_{j\neq i}(x_i - x_j)} ,
\]
de donde $c_{n-1} = \sum_i y_i/\prod_{j \neq i}(x_i - x_j)$ --- otra
vez la fórmula de las \omterm{pb:b1:vspaces:1}{diferencias divididas}.

\textbf{8.} $L_3 \leftarrow L_3 - L_1$ da las filas
$(1, x_1, x_1^2)$, $(0, 1, 2x_1)$,
$(0,\ x_2 - x_1,\ (x_2-x_1)(x_2+x_1))$; desarrollando por la primera
columna y sacando factor $(x_2 - x_1)$:
\[
(x_2 - x_1)\begin{vmatrix} 1 & 2x_1\\ 1 & x_2 + x_1
\end{vmatrix} = (x_2 - x_1)(x_2 - x_1) = (x_2 - x_1)^2 .
\]
No nulo para $x_1 \neq x_2$: el \omterm{def:b1:det:system}{sistema lineal} que expresa
$P(x_1) = u$, $P'(x_1) = v$, $P(x_2) = w$ sobre los coeficientes de
$P \in \R_2[X]$ es de Cramer --- la interpolación de Hermite con un
nodo doble está bien planteada.

\textbf{9.} $P = a + bX + cX^2$ con $a = P(0) = 1$,
$b = P'(0) = 0$, $a + b + c = P(1) = 2$: $c = 1$, luego
$P = 1 + X^2$, único. Coherencia: aquí $x_1 = 0$, $x_2 = 1$ y el
determinante de la pregunta 8 es $(1 - 0)^2 = 1 \neq 0$.

\textbf{10.} Cálculo directo:
\[
C_2 = \frac{1}{(a_1+b_1)(a_2+b_2)} - \frac{1}{(a_1+b_2)(a_2+b_1)}
= \frac{(a_1+b_2)(a_2+b_1) - (a_1+b_1)(a_2+b_2)}
{\prod_{i,j}(a_i+b_j)} ,
\]
y el numerador se desarrolla como
$a_1b_1 + a_2b_2 - a_1b_2 - a_2b_1 = (a_2 - a_1)(b_2 - b_1)$:
diferencias entre sumas.

\textbf{11.} Para $i < n$, la nueva entrada de la fila $i$ es
\[
\frac{1}{a_i + b_j} - \frac{1}{a_n + b_j}
= \frac{a_n - a_i}{(a_i + b_j)(a_n + b_j)} .
\]
Sáquese factor $(a_n - a_i)$ de cada fila $i < n$ y después
$\frac1{a_n + b_j}$ de cada columna $j$: lo que queda tiene entradas
$\frac1{a_i + b_j}$ en las filas $i < n$ y constante $1$ en la fila
$n$ --- la matriz $M$, con el prefactor anunciado.

\textbf{12.} Sobre $M$, para $j < n$ la operación
$C_j \leftarrow C_j - C_n$ convierte la fila $n$ en
$(0, \dots, 0, 1)$ y, en la fila $i < n$,
\[
\frac{1}{a_i + b_j} - \frac{1}{a_i + b_n}
= \frac{b_n - b_j}{(a_i + b_j)(a_i + b_n)} .
\]
Sáquese factor $(b_n - b_j)$ de cada columna $j < n$ y
$\frac1{a_i + b_n}$ de cada fila $i < n$, y desarróllese después por
la última fila (con signo $(-1)^{n+n} = +1$): el determinante
restante es $C_{n-1}$. Reuniendo los factores de las preguntas
11--12:
\[
C_n = \frac{\prod_{i<n}(a_n - a_i)\,\prod_{j<n}(b_n - b_j)}
{\prod_{j}(a_n + b_j)\,\prod_{i<n}(a_i + b_n)}\;C_{n-1},
\]
y la inducción (con base $C_1 = \frac1{a_1+b_1}$) ensambla
exactamente el doble alternante de Cauchy: los factores
$(a_j - a_i)(b_j - b_i)$ para todas las parejas, entre todas las
sumas $(a_i + b_j)$.

\textbf{13.} La fórmula se anula si y solo si algún $a_j = a_i$ o
$b_j = b_i$: la matriz de Cauchy es invertible si y solo si las dos
familias son distintas dos a dos. Hilbert: $a_i = i$, $b_j = j - 1$.
Para $n = 2$: numerador $(2-1)(1-0) = 1$, denominador
$1\cdot2\cdot2\cdot3 = 12$: $\det H_2 = \frac1{12}$. Para $n = 3$:
numerador $\bigl[(1)(2)(1)\bigr]^2 = 4$, denominador
$(1\cdot2\cdot3) (2\cdot3\cdot4)(3\cdot4\cdot5) = 6\cdot24\cdot60 =
8640$: $\det H_3 = \frac{4}{8640} = \frac1{2160}$. Inversa para
$n = 2$:
\[
H_2^{-1} = 12\begin{pmatrix} \frac13 & -\frac12\\[2pt]
-\frac12 & 1\end{pmatrix}
= \begin{pmatrix} 4 & -6\\ -6 & 12 \end{pmatrix},
\]
todas enteras (fenómeno cierto para todo $n$).

\textbf{14.} La matriz del sistema es la de Cauchy, invertible por la
pregunta 13 cuando los $b_j$ (y los $a_i$) son distintos dos a dos:
solución única. Interpretación: una función racional
$R = \sum_j \frac{c_j}{X + b_j}$ con \omterm{def:b1:fractions:field}{polos} simples queda determinada
por $n$ de sus valores $R(a_1), \dots, R(a_n)$ y, recíprocamente,
cualquier hoja de datos así se realiza exactamente una vez --- la
contrapartida por muestreo del teorema de existencia y unicidad de
las fracciones simples (\cref{thm:b1:fractions:complex}).

\textbf{15.} Vieta para $X^3 + pX + q$:
$\lambda_1 + \lambda_2 + \lambda_3 = 0$,
$\sum_{i<j}\lambda_i\lambda_j = p$, luego $p_1 = 0$ y
$p_2 = p_1^2 - 2p = -2p$. Cada raíz cumple
$\lambda^3 = -p\lambda - q$; sumando: $p_3 = -p\,p_1 - 3q = -3q$. Y
multiplicando por $\lambda$ y sumando: $p_4 = -p\,p_2 - q\,p_1 = 2p^2$.

\textbf{16.} Por la pregunta 5 (la identidad $S = W^{\mathsf T}W$ y
$\det S = V^2$ valen sobre $\C$),
\[
V^2 = \begin{vmatrix}
3 & 0 & -2p\\
0 & -2p & -3q\\
-2p & -3q & 2p^2
\end{vmatrix}
= 3\bigl(-4p^3 - 9q^2\bigr) + (-2p)\bigl(0 - 4p^2\bigr)
= -4p^3 - 27q^2 ,
\]
desarrollando por la primera fila.

\textbf{17.} Una raíz múltiple significa dos $\lambda_i$ iguales, es
decir, $V = 0$, es decir,
$\operatorname{disc} = -4p^3 - 27q^2 = 0$. Para $X^3 - 3X + 2$:
$4(-3)^3 + 27\cdot4 = -108 + 108 = 0$, acorde con la raíz doble $1$
de $(X-1)^2(X+2)$.

\textbf{18.} Las raíces no reales de una cúbica real vienen en
parejas \omterm{def:b1:complex:field}{conjugadas}, así que solo se dan dos casos cuando
$\operatorname{disc} \neq 0$. Tres raíces reales distintas: $V$ es
real y no nulo, luego $\operatorname{disc} = V^2 > 0$. Una raíz real
$\lambda_1$ y $\lambda_3 = \conj{\lambda_2} \notin \R$: entonces
\[
(\lambda_2 - \lambda_1)(\lambda_3 - \lambda_1) =
\abs{\lambda_2 - \lambda_1}^2 > 0,
\qquad
\lambda_3 - \lambda_2 = -2\iu\,\operatorname{Im}\lambda_2 \neq 0,
\]
luego $V$ es un número imaginario puro no nulo y
$\operatorname{disc} = V^2 < 0$. Los dos signos caracterizan los dos
casos.

\textbf{19.} Tómese $b_i = a_i$ en el doble alternante: el numerador
es $\prod_{i<j}(a_j - a_i)^2 > 0$ y el denominador,
$\prod_{i,j}(a_i + a_j) > 0$ (todas las entradas positivas): el
determinante es positivo. (En lenguaje posterior: el núcleo
$\frac1{x+y}$ es definido positivo.)

\textbf{20.} Por las preguntas 2--3, el determinante vale
$V(2,4,7)/(0!\,1!\,2!) = \frac{(4-2)(7-2)(7-4)}{2} =
\frac{30}{2} = 15$. Directamente, la matriz es
\[
\begin{pmatrix}
1 & 2 & 1\\
1 & 4 & 6\\
1 & 7 & 21
\end{pmatrix},
\qquad
\det = (84 - 42) - 2(21 - 6) + (7 - 4) = 42 - 30 + 3 = 15 .
\]

\textbf{21.} Supóngase $\sum_i c_i\,(\lambda_i^{\,k})_{k} = 0$ como
sucesión. Leyendo $k = 0, 1, \dots, n-1$ se obtiene
$W^{\mathsf T}c = 0$ con $W = (\lambda_i^{\,j-1})$ invertible
($\det = V \neq 0$, con los $\lambda_i$ distintos): $c = 0$. Las
sucesiones geométricas son \omterm{def:b1:vspaces:free}{libres}.

\textbf{22.} $a = b = (1, 2, 3)$: numerador
$\bigl[(2-1)(3-1) (3-2)\bigr]^2 = 4$; denominador
$\prod_{i,j}(i + j) = (2\cdot3\cdot4)(3\cdot4\cdot5)(4\cdot5\cdot6)
= 24\cdot60\cdot120 = 172800$. Por tanto,
$\det\bigl(\frac1{i+j}\bigr) = \frac{4}{172800} = \frac1{43200}$.

\textbf{23.} Si $x_i = x_j$, el intercambio de las dos variables deja
fijo el punto, pero tiene que cambiar el signo de $F$: $F = -F$,
luego $F = 0$ ahí. \omterm{def:b1:arith:divides}{Divisibilidad}: véase $F$ como \omterm{def:b1:poly:def}{polinomio} en la
única variable $x_n$ con coeficientes en las demás; se anula en los
$n - 1$ «valores» $x_1, \dots, x_{n-1}$, de modo que factorizar
repetidamente (\cref{thm:b1:poly:factor}) da
$F = \prod_{i<n}(x_n - x_i)\cdot G$ con $G$ \omterm{def:b1:poly:def}{polinomio}. El prefactor
es invariante por los intercambios de dos índices $i, j < n$, así que
$G$ es alternado en $x_1, \dots, x_{n-1}$, y la inducción remata:
$\prod_{i<j}(x_j - x_i)$ \omterm{def:b1:arith:divides}{divide} a $F$.

\textbf{24.} $D = \det(x_i^{\,j-1})$ es un \omterm{def:b1:poly:def}{polinomio} en los $x_i$;
intercambiar dos variables intercambia dos filas, luego $D$ es
alternado y, por la pregunta 23, $D = c\,\prod_{i<j}(x_j - x_i)$ para
cierto \omterm{def:b1:poly:def}{polinomio} $c$. Grados totales: $D$ tiene grado
$\leq 0 + 1 + \dots + (n-1) = \binom n2$ y el producto tiene grado
exactamente $\binom n2$: $c$ es una constante. El monomio
$x_2\,x_3^2\cdots x_n^{\,n-1}$ tiene coeficiente $1$ en $D$ (el
producto diagonal) y $1$ en el producto (elíjase en cada factor la
variable de índice mayor): $c = 1$, y la fórmula de Vandermonde cae
sin ninguna inducción.

\textbf{25.} (i) La alternancia ---el axioma «dos columnas iguales
matan el determinante»--- es el motor: produjo todos los factores
$(x_j - x_i)$, $(a_j - a_i)$, $(b_j - b_i)$ del problema. (ii) La
identidad $\det S = V^2$ sustituye las raíces complejas
individuales, inalcanzables, por sus sumas de potencias, que son
\omterm{def:b1:poly:def}{polinomios} reales en los coeficientes, de modo que la distinción
pasa a ser el signo de un número real calculable. (iii) El
\omterm{ex:b1:det:vandermonde}{determinante de Vandermonde} gobierna la interpolación
polinómica, y el \omterm{pb:b1:det:1}{determinante de Cauchy} gobierna las fracciones
simples y las funciones racionales muestreadas (con la
\omterm{pb:b1:det:1}{matriz de Hilbert} como su caso particular más famoso). (iv) Todo
\omterm{def:b1:poly:def}{polinomio} alternado es divisible por $\prod_{i<j}(x_j - x_i)$, y un
recuento de grados fija después un \omterm{def:b1:poly:def}{polinomio} así salvo una constante
--- y por eso este producto no deja de reaparecer allí donde un
determinante se anula sobre las coincidencias. El teorema de la
parte III es el \emph{doble alternante de Cauchy}.
\end{solution}
